We will now reformulate \eqref{eq:BM} to incorporate a multidimensional generalization of the ramp loss.
To do this, first observe that the objective \eqref{eq:BM-obj} can be written as
\begin{equation*}
    \sum_{i \in \cP} \zeta_i + \sum_{i \in \cN} \sum_{k \in \K} \gamma_{ik} w_k,
\end{equation*} where
\[\gamma_{ik} = \begin{cases}
                    1, & \text{if $k \in \K_i$,} \\
                    0, & \text{otherwise}.
\end{cases}\]
In $\gamma_{ik}$, we recognize a hard-margin loss.
Specifically, for a negative point $i \in \cN$, the value $\gamma_{ik}$ is exactly the loss induced if $k$ is used.
Similarly, constraint \eqref{eq:BM_orig-misP} can be written as
\begin{equation*}
    \zeta_i + \sum_{k \in \K} \gamma_{ik} w_k \geq 1, \quad \forall i \in \cP.
\end{equation*}
If a positive point $i \in \cP$ is uncovered, this implies the sum will be zero and hence $\zeta_i = 1$, which corresponds to an induced loss of $1$.
The value $\gamma_{ik}$ represents a reduction of this loss if rule $k$ is used.
If $k$ covers the point, this value is $1$.
Using rule $k$ then allows $\zeta_i$ to be reduced to $0$.

By defining $\bm{\gamma}$ suitably, we can introduce a multidimensional variant of the ramp loss.
When doing so, the values of $\bm{\gamma}$ need no longer be binary, hence $\bm{\zeta}$ should no longer be binary.
Furthermore, the ramp loss of a point is anywhere between $0$ and $2$.
The ramp loss model is therefore given by
\begin{subequations}
    \label{eq:RM}
    \begin{align}
        \mathbf{\min} \quad & \sum_{i \in \cP} \zeta_i + \sum_{i \in \cN} \sum_{k \in \K} \gamma_{ik} w_k \label{eq:RM-obj} \\
        \mathbf{\text{s.t.}} \quad & \zeta_i + \sum_{k \in \K} \gamma_{ik} w_k \geq 2, & \forall i \in \cP \label{eq:RM-misP} \\
        & \sum_{k \in \K} c_k w_k \leq C, & \label{eq:RM-complex} \\
        & w \in \{0, 1\}^{\K}, \zeta \in [0,2]^{\cP}, \label{eq:RM-vars}
    \end{align}
\end{subequations}
To define $\bm{\gamma}$, we first define a decision function $\bm{d}: \P(\L) \to \R^\I$ such that $d_i(k)$ is
non-negative if rule $k$ covers point $i$ and negative otherwise.
Furthermore, the value $d_i(k)$ quantifies how far $i$ is from the boundary of $k$.
Define
\begin{equation}\label{eq:decision_function}
    d_i(k) = \umin{(j,\alpha) \in k} X_{ij} - \alpha, \tquad i \in \I, \; k \in \K\setminus\{\emptyset\}.
\end{equation}
If $k$ covers $i$, then any element in the minimum is non-negative, meaning $d_i(k)$ is the distance to the closest threshold.
If rule $k$ does not cover $i$, then there exists some clause $(j,\alpha) \in k$ that cuts off $i$.
Choose such $(j,\alpha)$.
It follows that $X_{ij} - \alpha < 0$ and so $d_i(k) < 0$.
By extension, we have in this case that
\[
    d_i(k) = -\umax{(j,\alpha) \in k \fw X_{ij} < \alpha} \alpha - X_{ij}, \tquad i \in \I.
\]
This means that $d_i(k)$ is the (negated) distance to the furthest unsatisfied threshold.
Therefore, the value $-d_i(k)$ can be interpreted as a distance to the covered region.
To account for the empty rule, we set
\begin{equation}
    d_i(\emptyset) = \infty, \tquad i \in \I. \label{eq:decision_function_empty_rule}
\end{equation}
Denote by $\beta$ the margin width, as illustrated in \autoref{fig:loss_functions}b.
While the margin width is a variable in the Support Vector Machine~\cite{belotti_spatial_2024},
it is a constant in our model.
We define
\begin{equation}\label{eq:gamma_ik}
    \gamma_{ik} = \begin{cases}
                      0, & \text{if $d_i(k) \leq -\beta$}, \\
                      1 + \frac{1}{\beta}d_i(k), & \text{if $d_{i}(k) \in (-\beta, \beta]$}, \\
                      2, & \text{if $d_i(k) > \beta$},
    \end{cases} \tquad i \in \I, \; k \in \K.
\end{equation}
This definition realizes a generalized ramp loss in the master problem \eqref{eq:RM}.

We refer to this model as the ``additive'' ramp loss model.
This name is due to the additive nature of loss reduction for positive points.
Recall \eqref{eq:RM-misP}.
Several rules, when used together, can all contribute to reducing a $\zeta$-variable.
For example, if point $i$ lies in several margins, the corresponding loss $\zeta_i$ can still be $0$,
which is equivalent to being far inside the covered region of the ruleset.


As $\beta \to 0,$ the ramp loss converges pointwise to the hard-margin loss (multiplied by $2$)
for all values of $X_{ij}$ except the corresponding $\alpha$.
Assuming no feature values lies on such threshold,
it follows that the ramp loss model is equivalent to the threshold model when $\beta$ is arbitrarily small.
To see this, observe that for $\beta$ small enough, $\bm{\gamma}$ and $\bm{\zeta}$ are those used in
the threshold model, multiplied by $2$.

In a realistic setting, it is to be expected that few, or none, of the datapoints lie exactly on a threshold.
There are two reasons for this.
First, continuous measurements typically include a measurement error,
making it less likely for several datapoints to have an equal feature value.
Second, the linear interpolation used in the quantile function is likely to yield a non-observed value.
Therefore, in a practical setting with continuous features,
the ramp loss model with $\beta$ arbitrarily small is practically equivalent to the threshold model.
For convenience, we define for the special case $\beta = 0$,
\begin{equation}\label{eq:ramp_gamma_mw_0}
    \begin{split}
        \gamma_{ik} &= 2 \cdot \Ind(d_i(k) \geq 0) \\
        &=
        \begin{cases}
            0, & \text{if $\exists (j,\alpha)\in k \fw X_{ij} < \alpha$},\\
            2, & \text{if $X_{ij} \geq \alpha \quad \forall (j,\alpha) \in k$},\\
        \end{cases} \tquad i \in \I, \; k \in \K.
    \end{split}
\end{equation}
For this choice of $\bm{\gamma}$, the model is exactly equivalent to the threshold model.
We then use the implementation of the ramp loss model with $\beta = 0$ to solve the threshold model.
The goal is now to express the pricing problem in terms of $\bm{\gamma}$ and study how this problem may be solved.